\documentstyle[% 


righttag% 


,11pt% 


,amssymb% 


,russian%               remove in Eng. 


,izvru%                 Set izven% in Eng. 


% ,izvru-ks             for brief communications 


]{amsart} 


 


%  


 


\newcommand{\tg}{\operatorname{tg}} 


\newcommand{\tts}[1]{} 


 


\theoremstyle{definition} 


\theorembodyfont{\rm} 


\newtheorem{cond}{\hskip\parindent Условие} 


\newtheorem{notenonum}{\hskip\parindent Замечание} 


\renewcommand{\thenotenonum}{} 


 


%\hoffset=-15mm%                  For Eng. 


\setcounter{page}{72} 


\god{2005} 


\nomer{10~(521)} %                        Remove in Eng. 


%\nomer{Vol.\,49, No.\,10}%               Set in Eng. 


%\str{\pageref{begin}--\pageref{end}}%  Set in Eng. 


\udk{517.956} 


 


\author{\it Р.С.\,Хайруллин} 


% NB:   ^^ remove in Eng. 


\title{О задаче типа Геллерстедта 


для уравнения второго рода} 


 


\date{} 


%\pagestyle{myheadings}%         Set in Eng. 


\pagestyle{plain} %              Remove in Eng. 


\begin{document} 


%\thispagestyle{plain}%          Set in Eng. 


\maketitle 


\label{begin} 


 


 


 


\section{Постановка задачи} 


 


Рассмотрим уравнение 


\begin{equation} 


u_{xx}+yu_{yy}+\alpha u_y=0,\quad \alpha\le -1/2, 


\end{equation} 


в смешанной области $D$, эллиптическая часть которой $D_1$ совпадает со 


всей верхней полуплоскостью, а гиперболическая часть состоит из двух 


бесконечных треугольников: $D_2$, ограниченного характеристиками $y=0$ и 


$x-2\sqrt{-y}=0$, и $D_3$, ограниченного характеристиками $y=0$ и 


$x+2\sqrt{-y}=0$. Задачи Геллерстедта для уравнения (1) в случае 


ограниченных областей рассматривались в работах [1]--[3]. В данной статье 


исследуется случай неограниченной области, и в ней продолжены 


исследования, начатые в работах [4], [5], в которых были решены 


аналогичные задачи Трикоми. 


 


Через $n$ и $m$ обозначим натуральные числа, удовлетворяющие неравенствам 


$-1/2<\alpha+n=\alpha_0\le1/2$, $0<2\alpha+m-1=\delta\le1$. Очевидно, 


$m=2n+2$, $\delta=2\alpha_0+1$ при $-1/2<\alpha_0\le0$ и $m=2n+1$,


$\delta=2\alpha_0$ при $0<\alpha_0\le 1/2$. 


\medskip 


 


{\bf Задача $G_\alpha$.} {\it В области $D$ найти функцию $u(x,y)$ со 


свойствами 


\begin{itemize} 


\item[1)] $u(x,y)$ принадлежит $C(D\cup\{x^2+4y=0\});$ 


\item[2)] имеют место соотношения 


\begin{equation} 


u=o(R^{2-2\alpha}),\quad u_x=o(R^{1-2\alpha}),\quad 


u_y=o(R^{-2\alpha}) 


\end{equation} 


при $R\to+\infty$, где $R^2=x^2+4y$, $(x,y)\in D_1;$ 


\item[3)] $u(x,y)$ принадлежит $C^2(D_1\cup D_2\cup D_3)$ и удовлетворяет 


уравнению $(1)$ в $D_1\cup D_2\cup D_3;$ 


\item[4)] существуют пределы $(x\ne0)$ 


\begin{equation} 


\nu_i(x)=\lim_{y\to0,\ (x,y)\in D_i}|y|^\alpha 


[u(x,y)-A_\alpha(x,y,\tau)]_y, 


\end{equation} 


$i=1,2,3$, и выполняются условия склеивания 


\begin{alignat}{2} 


\nu_1(x)&=(-1)^n \nu_2(x), &\quad x&>0,\\ 


\nu_1(x)&=(-1)^n\nu_3(x), &\quad x&<0, 


\end{alignat} 


где 


\begin{gather} 


\tau(x)=u(x,0),\\ 


A_\alpha(x,y,\tau)=\sum^{[m/2]}_{s=1} 


\frac{\tau^{(2s)}(x)(-1)^s}{(\alpha)_s s!}y^s,\quad 


\alpha\ne -n,\notag\\ 


A_\alpha(x,y,\tau)=\sum^n_{s=1} 


\frac{\tau^{(2s)}(x)(-1)^s}{(\alpha)_s s!}y^s- 


\frac{\tau^{(2n+2)}(x)}{n!(n+1)!}y^{n+1}\bigg[\ln|y| 


-\sum^{n+1}_{s=1}\frac{1}{s}\bigg],\quad \alpha=-n,\notag 


\end{gather} 


$[\,\cdot\,]$ --- целая часть числа, $(\alpha)_0=1$, 


$(\alpha)_s=\alpha(\alpha+1)(\alpha+2)\dotsb(\alpha+s-1);$ 


\item[5)] $u(x,y)$ удовлетворяет краевому условию 


$$ 


u(x,y)|_{x^2+4y=0}=\psi(x), 


$$ 


где $\psi(x)$ --- заданная функция. 


\end{itemize}


} 


\medskip 


 


Предполагаем, что удовлетворяется 


 


 


\begin{cond} 


Функция $\psi(x)$ принадлежит классу $C^n(-\infty,+\infty)\cup 


C^{n+1,\lambda}(-\infty,0)\cup C^{n+1,\lambda}(0,+\infty)$, 


$\lambda>1/2-\alpha_0$, производная $\psi^{(n+1)}(x)$ может иметь 


особенность при $x=0$ порядка ниже $\alpha_0+1/2$, если 


$-1/2<\alpha_0\le0$, и ниже $\delta$, если $0<\alpha_0\le1/2$, и должна 


иметь представление $\psi^{(n+1)}(x)=O(|x|^{-\beta})$ при $x\to\infty$, 


где $\beta>\alpha_0$. 


\end{cond} 


 


На функции $\tau(x)$ и $\nu_i(x)$ наложим следующие условия 2 и 3. 


 


 


\begin{cond} 


Функция $\tau(x)$ принадлежит классу $C^n(-\infty,+\infty)\cup 


C^{m,\gamma}(-\infty,0)\cup C^{m,\gamma}(0,+\infty)$, $\gamma>1-\delta$, 


производная $\tau^{(n+1)}(x)$ может иметь особенность при $x=0$ порядка 


ниже $\alpha_0+1/2$, если $-1/2<\alpha_0\le0$, и ниже $\delta$, если 


$0<\alpha_0\le1/2$, и должна иметь представление 


$\tau^{(n+1)}(x)=O(|x|^{-\beta})$ при $x\to\infty$. 


\end{cond} 


 


 


\begin{cond} 


Функции $\nu_i(x)$ непрерывны в соответствующих областях 


определения и могут иметь особенность при $x=0$ порядка ниже 


$3/2-\alpha$. 


\end{cond} 


 


 


\section{Основные соотношения между $\tau$ и $\nu$} 


 


Задачу будем решать методом интегральных уравнений, поэтому понадобятся 


основные соотношения между $\tau$ и $\nu$ из каждой подобласти. Для 


вывода этих соотношений рассмотрим вспомогательные задачи. В 


эллиптической подобласти используем решение задачи Дирихле с краевым 


условием (6). Решение ищется в классе функций, удовлетворяющих на 


бесконечности соотношениям (2). Указанная задача имеет единственное 


решение 


$$ 


u(x,y)=\frac{\Gamma(3/2-\alpha)y^{1-\alpha}} 


{\sqrt{\pi}\Gamma(1-\alpha)4^{\alpha-1}} 


\int^{+\infty}_{-\infty}\tau(\xi) 


[(x-\xi)^2+4y]^{\alpha-3/2}d\xi.


$$ 


В классе ограниченных функций это решение построено в [6], [7]. Оно 


остается справедливым и для неограниченных функций из требуемого класса. 


В этом можно убедиться, используя метод функции Грина. Аналогичная задача 


для области, совпадающей с первым квадрантом, решена в [4]. 


 


Перейдем к гиперболическим подобластям. Здесь используем решение 


задачи типа Коши с начальными условиями (6), (3). Оно имеет вид [8] 


$$ 


u(x,y)=B_\alpha(x,y,\tau)- 


\frac{2\Gamma(2-2\alpha)}{\Gamma^2(3/2-\alpha)}(-y)^{1-\alpha} 


\int^1_0\nu_i(\zeta)[\xi(1-\xi)]^{1/2-\alpha}d\xi, 


$$ 


где $(x,y)\in D_i$, $i=2,3$, $\zeta=x-2\sqrt{-y}(1-2\xi)$, 


\begin{gather*} 


\displaybreak[3]


B_\alpha(x,y,\tau)=\frac{1}{2}\sum^n_{s=0} 


\frac{n!(2n-s)!2^{2s}}{s!(n-s)!(2n)!}(-y)^{s/2} 


[\tau^{(s)}(x-2\sqrt{-y})+(-1)^s\tau^{(s)} 


(x+2\sqrt{-y})],\ \ \alpha_0=1/2,\\ 


% 


B_\alpha(x,y,\tau)=\sum^n_{s=0} 


\frac{C^s_{n+1}\Gamma(\alpha_0+1)4^{\alpha_0+s}} 


{\sqrt{\pi}(\alpha)_s\Gamma(\alpha_0+s+1/2)}(-y)^s\int^1_0 


\tau^{(2s)}(\zeta)[\xi(1-\xi)]^{\alpha_0+s-1/2}d\xi+\\ 


% 


+\frac{\Gamma(\alpha)4^{\alpha+m-1}} 


{\sqrt{\pi}\Gamma(\alpha+m-1/2)}(-y)^{m/2}\int^1_0 


\tau^{(m)}(\zeta)[\xi(1-\xi)]^{\alpha+m-3/2} 


[\xi+(-1)^m(1-\xi)]d\xi,\ \ \delta<1,\\ 


% 


B_\alpha(x,y,\tau)=\sum^n_{s=0} 


\frac{C^s_{n+1}4^s}{\sqrt{\pi}(\alpha)_s\Gamma(s+1/2)} 


(-y)^s\int^1_0 \tau^{(2s)}(\zeta)[\xi(1-\xi)]^{s-1/2}d\xi+\\ 


% 


+\frac{(-1)^n2^{2n+3}(-y)^{n+1}} 


{\sqrt{\pi}\Gamma(n+3/2)n!}\int^1_0\tau^{(2n+2)} 


(\zeta)[\xi(1-\xi)]^{n+1/2}\bigg[\ln 16\sqrt{-y}\xi(1-\xi)- 


\sum^{n+1}_{s=1}\frac{2}{2s-1}\bigg]d\xi,\ \ \alpha_0=0. 


\end{gather*} 


Из этих решений выводятся требуемые соотношения. Их вид 


определяется следующими леммами. Они доказываются по аналогии 


с соответствующими утверждениями работ [4], [5]. 


 


 


\begin{lem} 


При $x>0$ основное соотношение из эллиптической подобласти 


имеет вид\linebreak $(a>x)$ 


\begin{gather*}


\Gamma(1-\alpha)\nu_1(x)=k\bigg[\Gamma(2\alpha-2) 


\frac{d^{m-n}}{dx^{m-n}}(D^{-\delta}_{0x}\tau^{(n+1)}(x) 


+(-1)^{m+1}D^{-\delta}_{xa}\tau^{(n+1)}(x))+\\ 


% 


+\frac{1}{(2-\delta)_m}\frac{d^{m-n-1}}{dx^{m-n-1}}\bigg( 


\int^{+\infty}_a \tau^{(n+1)}(\xi)(\xi-x)^{\delta-2}d\xi+ 


(-1)^m\int^0_{-\infty}\tau^{(n+1)}(\xi)(x-\xi)^{\delta-2} 


d\xi\bigg)\bigg],\ \ \delta<1,\\ 


% 


\Gamma(1-\alpha)\nu_1(x)=\frac{k}{m!}\bigg[ 


\frac{d^{m-n}}{dx^{m-n}}\bigg((-1)^m\int^x_0 


\tau^{(n+1)}(\xi)\ln(x-\xi)d\xi-\\ 


% 


-\int^a_x \tau^{(n+1)}(\xi)\ln(\xi-x)d\xi\bigg)+ 


\frac{d^{m-n-1}}{dx^{m-n-1}}\bigg(\int^{+\infty}_a 


\frac{\tau^{(n+1)}(\xi)d\xi}{\xi-x}-(-1)^m\int^0_{-\infty} 


\frac{\tau^{(n+1)}(\xi)d\xi}{\xi-x}\bigg)\bigg],\ \ \delta=1, 


\end{gather*} 


где $k=\Gamma(3/2-\alpha)(1-\alpha)/\sqrt{\pi}4^{\alpha-1}$. 


\end{lem} 


 


 


\begin{lem} 


При $x<0$ основное соотношение из эллиптической подобласти имеет вид \linebreak


$(-a<x)$ 


\begin{gather*} 


\Gamma(1-\alpha)\nu_1(x)=k\bigg[\Gamma(2\alpha-2) 


\frac{d^{m-n}}{dx^{m-n}}(D^{-\delta}_{-ax}\tau^{(n+1)}(x) 


+(-1)^{m+1}D^{-\delta}_{x0}\tau^{(n+1)}(x))+\\ 


% 


+\frac{1}{(2-\delta)_m}\frac{d^{m-n-1}}{dx^{m-n-1}}\bigg( 


\int^{+\infty}_0 \tau^{(n+1)}(\xi)(\xi-x)^{\delta-2}d\xi+ 


(-1)^m\int^{-a}_{-\infty}\tau^{(n+1)}(\xi)(x-\xi)^{\delta-2} 


d\xi\bigg)\bigg],\ \ \delta<1,\\ 


% 


\Gamma(1-\alpha)\nu_1(x)=\frac{k}{m!}\bigg[ 


\frac{d^{m-n}}{dx^{m-n}}\bigg((-1)^m\int^x_{-a} 


\tau^{(n+1)}(\xi)\ln(x-\xi)d\xi-\\ 


% 


-\int^0_x \tau^{(n+1)}(\xi)\ln(\xi-x)d\xi\bigg)+ 


\frac{d^{m-n-1}}{dx^{m-n-1}}\bigg(\int^{+\infty}_0 


\frac{\tau^{(n+1)}(\xi)d\xi}{\xi-x}-(-1)^m\int^{-a}_{-\infty} 


\frac{\tau^{(n+1)}(\xi)d\xi}{\xi-x}\bigg)\bigg],\ \ \delta=1.


\end{gather*} 


\end{lem} 


 


 


\begin{lem} 


Основное соотношение из гиперболической подобласти $D_2$ 


имеет вид 


\begin{align*} 


\Gamma(1-\alpha)\nu_2(x)&=\Gamma(\alpha) 


\frac{d^{m-n}}{dx^{m-n}}D^{-\delta}_{0x}\tau^{(n+1)}(x) 


-\Psi_\alpha(x,\psi^{(n+1)}),\ \ \alpha_0\ne0,\\ 


% 


\Gamma(1-\alpha)\nu_2(x)&=\frac{2(-1)^n}{n!} 


\frac{d^{m-n}}{dx^{m-n}}\int^x_0 \tau^{(n+1)}(\xi) 


\ln(x-\xi)d\xi-\Psi_\alpha(x,\psi^{(n+1)}),\ \ \alpha_0=0, 


\end{align*} 


$$ 


\tau^{(s)}(0)=\psi^{(s)}(0)2^{-s}(2\alpha-1)_s/(\alpha-1/2)_s, 


\ \ s=\ov{0,n}, 


$$ 


где 


$$ 


\Psi_\alpha(x,\psi^{(n+1)})=2^{1-2\alpha-n} 


\sqrt{\pi}x^{\alpha-1/2}D^{1/2-\alpha_0}_{0x}\psi^{(n+1)}(x/2). 


$$ 


\end{lem} 


 


 


\begin{lem} 


Основное соотношение из гиперболической подобласти $D_3$ 


имеет вид 


\begin{align*} 


\Gamma(1-\alpha)\nu_2(x)&=\Gamma(\alpha)(-1)^{m+1} 


\frac{d^{m-n}}{dx^{m-n}}D^{-\delta}_{x0}\tau^{(n+1)}(x) 


-\Psi_\alpha(x,\psi^{(n+1)}),\ \ \alpha_0\ne0,\\ 


% 


\Gamma(1-\alpha)\nu_3(x)&=\frac{2(-1)^{n+1}}{n!} 


\frac{d^{m-n}}{dx^{m-n}}\int^0_x \tau^{(n+1)}(\xi) 


\ln(\xi-x)d\xi-\Psi_\alpha(x,\psi^{(n+1)}),\ \ \alpha_0=0, 


\end{align*} 


$$ 


\tau^{(s)}(0)=\psi^{(s)}(0)2^{-s}(2\alpha-1)_s/(\alpha-1/2)_s, 


\ \ s=\ov{0,n}, 


$$ 


где 


$$ 


\Psi_\alpha(x,\psi^{(n+1)})=2^{1-2\alpha-n} 


\sqrt{\pi}(-x)^{\alpha-1/2}(-1)^{n+1} 


D^{1/2-\alpha_0}_{x0}\psi^{(n+1)}(x/2). 


$$ 


\end{lem} 


 


Здесь $D^l_{bc}$ --- оператор дробного дифференцирования по 


Риману--Лиувиллю 


при $l\ge0$ и дробного интегрирования при $l<0$ (напр.,[9]). 


 


 


\begin{lem} 


Если $\psi(x)$ удовлетворяет условию $1$, то функция 


$\Psi_\alpha(x,\psi^{(n+1)})$ может иметь при $x=0$ особенность порядка 


ниже $\frac32-\alpha$, если $-\frac12<\alpha_0\le0$, и ниже $n+1$, если 


$0<\alpha_0\le\frac12$, а при $x\to\infty$ имеет нуль порядка выше 


$1-\alpha$. 


\end{lem} 


 


 


\section{Вывод интегральных уравнений и решение задачи $G_\alpha$} 


 


Приступим к выводу интегральных уравнений и их решению. 


 


 


\begin{thm} 


Если функция $\psi(x)$ удовлетворяет условию $1$, то решение 


задачи $G_\alpha$ в классе функций, удовлетворяющих условиям $2$, $3$, 


редуцируется к решению следующей системы интегральных уравнений\/${:}$


\begin{align} 


\mu_1(x)\tg\frac{\pi\alpha_0}{2}-\frac{1}{\pi} 


\int^{+\infty}_0\frac{\mu_1(\xi)d\xi}{\xi-x}- 


\frac{(-1)^m}{\pi}\int^{+\infty}_0 


\frac{\mu_2(\xi)d\xi}{\xi+x}&=g_1(x) 


+\sum^{m-n-2}_{s=0}c^1_s x^s,\\ 


% 


\mu_2(x)\tg\frac{\pi\alpha_0}{2}-\frac{1}{\pi} 


\int^{+\infty}_0\frac{\mu_2(\xi)d\xi}{\xi-x}- 


\frac{(-1)^m}{\pi}\int^{+\infty}_0 


\frac{\mu_1(\xi)d\xi}{\xi+x}&=g_2(x) 


+\sum^{m-n-2}_{s=0}c^2_s x^s


\end{align} 


с дополнительными условиями 


\begin{equation} 


\tau^{(s)}(0)=a_s,\quad s=\ov{0,n}, 


\end{equation} 


где 


\begin{alignat}{2} 


\mu_1(x)&=\tau^{(n+1)}(x)x^{\delta-1},&\quad x&>0,\\ 


\mu_2(-x)&=\tau^{(n+1)}(x)(-x)^{\delta-1},&\quad x&<0, 


\end{alignat} 


% 


\begin{align*} 


g_1(x)&=\frac{(-1)^{m+1}\Gamma(1-\alpha)x^{\delta-1}} 


{\pi(m-n-2)!}D^{\delta-1}_{0x}\int^{+\infty}_x 


\Psi_\alpha(\xi,\psi^{(n+1)})(\xi-x)^{m-n-2}d\xi,\\ 


% 


g_2(-x)&=\frac{(-1)^{m-n}\Gamma(1-\alpha)(-x)^{\delta-1}} 


{\pi(m-n-2)!}D^{\delta-1}_{x0}\int_{-\infty}^x 


\Psi_\alpha(\xi,\psi^{(n+1)})(x-\xi)^{m-n-2}d\xi,


\end{align*} 


$$ 


a_s=\psi^{(s)}(0)2^{-s}(2\alpha-1)_s/(\alpha-1/2)_s,\quad 


s=\ov{0,n}, 


$$ 


$c^1_s$, $c^2_s$, $s=\ov{0,m-n-2}$, --- произвольные постоянные. 


\end{thm} 


 


При доказательстве используются леммы 1--4 и условия склеивания (4), (5). 


 


 


\begin{lem} 


Если функция $\psi(x)$ удовлетворяет условию $1$, то функции $g_i(x)$ 


могут иметь особенность при $x=0$ порядка ниже $1/2-\alpha_0$, если 


$-1/2<\alpha_0\le0$, и ниже $1$, если $0<\alpha_0\le1/2$, и имеют нуль 


при $x\to+\infty$ порядка выше $-\alpha_0$, если $-1/2<\alpha_0\le0$, и 


выше $1-\alpha_0$, если $0<\alpha_0\le1/2$. 


\end{lem} 


 


 


\begin{notenonum} 


Из условия 2 и соотношений (10), (11) следует, что функции 


$\mu_i(x)$ при $x=0$ и $x\to+\infty$ должны удовлетворять условиям, 


аналогичным приведенным в утверждении леммы~6. 


\end{notenonum} 


 


 


\begin{lem} 


При условиях леммы $6$ система уравнений $(7)$, $(8)$ имеет 


единственное решение из требуемого класса. 


\end{lem} 


 


 


\begin{pf} 


Умножив уравнение (8) на $(-1)^m$, сложим его с уравнением 


(7) и вычтем из него. В результате имеем 


\begin{align*} 


\rho_1(x)\tg\frac{\pi\alpha_0}{2}-\frac{1}{\pi} 


\int^{+\infty}_0\bigg[\frac{1}{\xi-x}+\frac{1}{\xi+x}\bigg] 


\rho_1(\xi)d\xi&=h_1(x)+\sum^{m-n-2}_{s=0} d^1_s x^s,\\ 


% 


\rho_2(x)\tg\frac{\pi\alpha_0}{2}-\frac{1}{\pi} 


\int^{+\infty}_0\bigg[\frac{1}{\xi-x}-\frac{1}{\xi+x}\bigg] 


\rho_2(\xi)d\xi&=h_2(x)+\sum^{m-n-2}_{s=0} d^2_s x^s,


\end{align*} 


где 


\begin{alignat*}{2} 


\rho_1(x)&=\mu_1(x)+\mu_2(x)(-1)^m,&\quad 


\rho_2(x)&=\mu_1(x)-\mu_2(x)(-1)^m,\\ 


h_1(x)&=g_1(x)+g_2(x)(-1)^m,&\quad 


h_2(x)&=g_1(x)-g_2(x)(-1)^m. 


\end{alignat*} 


Отсюда после замены $\xi^2=\sigma$, $x^2=\eta$, 


$$ 


\rho_1(x)=v_1(\eta),\ \ h_1(x)=b_1(\eta), 


\ \ \rho_2(x)/x=v_2(\eta),\ \ h_2(x)/x=b_2(\eta) 


$$ 


получим 


\begin{align} 


v_1(\eta)\tg\frac{\pi\alpha_0}{2}-\frac{1}{\pi} 


\int^{+\infty}_0\frac{v_1(\sigma)d\sigma}{\sigma-\eta} 


&=b_1(\eta)+\sum^{m-n-2}_{s=0}d^1_s\eta^{s/2},\\ 


% 


v_2(\eta)\tg\frac{\pi\alpha_0}{2}-\frac{1}{\pi} 


\int^{+\infty}_0\frac{v_2(\sigma)d\sigma}{\sigma-\eta} 


&=b_2(\eta)+\sum^{m-n-2}_{s=0}d^2_s\eta^{(s-1)/2}.


\end{align} 


 


Для существования решений уравнений вида (12), (13) необходимо 


обращение в нуль правых частей при $x\to+\infty$ [10], поэтому имеем 


$d^1_s=0$, $s=\ov{0,m-n-2}$, 


и $d^2_s=0$, $s=\ov{1,m-n-2}$. Выполним подстановку [10] 


\begin{alignat*}{2} 


\eta&=\zeta/(1-\zeta),&\quad \sigma&=t/(1-t),\\ 


v_i(\eta)&=z_i(\zeta)(1-\zeta),&\quad 


b_i(\eta)&=w_i(\zeta)(1-\zeta) 


\end{alignat*} 


и получим 


\begin{align} 


z_1(\zeta)\tg\frac{\pi\alpha_0}{2}-\frac{1}{\pi}\int^1_0 


\frac{z_1(t)dt}{t-\zeta}&=w_1(\zeta), \\ 


% 


z_2(\zeta)\tg\frac{\pi\alpha_0}{2}-\frac{1}{\pi}\int^1_0 


\frac{z_2(t)dt}{t-\zeta}&=w_2(\zeta)+ 


\frac{d^2_0}{\sqrt{\zeta(1-\zeta)}}.


\end{align} 


Функция $w_1(\zeta)$ может иметь особенности при $\zeta=0$ порядка ниже 


$(1-2\alpha_0)/4$, если $-1/2<\alpha_0\le0$, и ниже $1/2$, если 


$0<\alpha_0\le1/2$, при $\zeta=1$ --- порядка ниже $1+\alpha_0/2$, если 


$-1/2<\alpha_0\le0$, и ниже $(1+\alpha_0)/2$, если $0<\alpha_0\le1/2$. 


Функция $w_2(\zeta)$ может иметь особенности при $\zeta=0$ порядка ниже 


$(3-2\alpha_0)/4$, если $-1/2<\alpha_0\le0$, и ниже~1, если 


$0<\alpha_0\le1/2$, при $\zeta=1$ --- порядка ниже $(1+\alpha_0)/2$, если 


$-1/2<\alpha_0\le0$, и ниже $\alpha_0/2$, если $0<\alpha_0\le1/2$. 


Аналогичные особенности допускаются у функций $z_1(\zeta)$ и 


$z_2(\zeta)$. Из того, что у функций $z_2(\zeta)$ и $w_2(\zeta)$ при 


$\zeta=1$ допускаются особенности порядка ниже $1/2$, следует равенство 


$d^2_0=0$. В требуемых классах уравнения (14), (15) имеют единственные 


решения. Эти решения записываются для уравнения (15) по формуле решения, 


ограниченного при $\zeta=0$. По этой же формуле оно записывается и для 


уравнения (14) в случае $0<\alpha_0\le1/2$. Если же $-1/2<\alpha_0\le0$, 


то для уравнения (14) решение записывается по формуле решения, 


ограниченного при $\zeta=0$. 


\end{pf} 


 


С помощью леммы 7 нетрудно доказывается 


 


 


\begin{thm} 


При условиях теоремы $1$ задача $(7)$--\/$(9)$ имеет 


единственное решение, удовлетворяющее условию $2$. 


\end{thm} 


 


Итогом является 


 


 


\begin{thm} 


Если функция $\psi(x)$ удовлетворяет условию $1$, то задача $G_\alpha$ 


имеет единственное решение, удовлетворяющее условиям $2$, $3$. 


\end{thm} 


 


Доказательство состоит в последовательном использовании теорем 1, 2. 


Единственность решения следует из однозначности определения функций 


$\tau(x)$ и $\nu_i(x)$ и единственности решений вспомогательных задач. 


 


 


\begin{thebibliography}{99} 


 


\bibitem{1} 


Хе Кан Чер. {\em О задаче Геллерстедта для одного уравнения смешанного 


типа} // Динамика сплошной среды. -- Новосибирск, 1976. 


-- Вып.\,26. -- С.\,134--141. 


\bibitem{2} 


Исамухамедов С.С. {\em Краевые задачи Геллерстедта для одного 


уравнения смешанного типа второго рода} // Дифференц.


уравнения с частн. произв. и их приложения. -- Ташкент, 


1977. -- С.\,33--40. 


\bibitem{3} 


Емелина И.Д. {\em Задачи типа Геллерстедта для одного однородного 


уравнения смешанного типа второго рода} // Тр. семин. по краев. 


задачам. --- Казань, 1983. -- Вып.\,20. -- С.\,93--103. 


\bibitem{4} 


Хайруллин Р.С. {\em К задаче Трикоми для уравнения смешанного типа 


второго рода} // Сиб. матем. журн. -- 1994. -- Т.35. -- \No\,4. -- 


С.\,927--936. 


\bibitem{5} 


Хайруллин Р.С. {\em Задача Трикоми для уравнения смешанного типа 


второго рода в случае неограниченной области} // Дифференц.


уравнения. -- 1994. -- Т.\,30. -- \No\,11. -- С.\,2010--2017. 


\bibitem{6} 


Джаиани Г.В. {\em Уравнение Эйлера--Пуассона--Дарбу}.


-- Тбилиси: Изд-во Тбилисск. ун-та, 1984. -- 80~с. 


\bibitem{7} 


Маричев О.И. {\em Сингулярные краевые задачи для уравнения 


Эйлера--Пуассона--Дарбу и двуосесимметричного уравнения 


Гельмгольца} // Тр. Всесоюзн. конф. по уравн. в частн. 


произв. -- М., 1978. -- С.\,373--374. 


\bibitem{8} 


Хайруллин Р.С. {\em Задача Трикоми для уравнения смешанного типа 


второго рода в случае нормальной области} // Дифференц.


уравнения. -- 1990. -- Т.\,26. -- \No\,8. -- С.\,1396--1407. 


\bibitem{9} 


Смирнов М.М. {\em Уравнения смешанного типа}.


-- М.: Высш. шк., 1985. -- 305 с. 


\bibitem{10} 


Гахов Ф.Д. {\em Краевые задачи}. -- М.: Наука, 1977. -- 640 с. 


 


\end{thebibliography} 


 


\vskip 2cm 


 


\postup{\it Казанская государственная}{\qquad \it Поступили} 


\postup{\it архитектурно-строительная}{\it первый вариант $08.05.2003$} 


\postup{\it академия}{\it окончательный вариант $19.02.2004$} 


 


\label{end} 


\end{document} 


